\begin{proposition}[The anchor-seed counterrange has only three
report-free petals]
\label{prop:double-cycle-anchor-seed-obstruction}
Fix an even $n\geq6$ in the double-cycle graph
\cref{eq:double-cycle-feed-edges}.  Run the exact canonical all-violations
gate from the anchor seed $s=e_{u_j}$ in the nonempty range
\begin{equation*}
 0<\rho<\frac{1-\alpha}{8}.
\end{equation*}
Before the first batch that intersects $W$, at most the three sites
$f_j,f_{j-1},f_{j+1}$ can be admitted.  More sharply, if a canonical batch
first intersects $\{f_{j-2},f_{j+2}\}$, then that same batch contains $w_j$
unless a report was admitted earlier.  Consequently this anchor-seed
counterexecution has no report-free epoch with $\Omega(n)$ distinct petal
admissions.  In particular it does not rescue the double-cycle feed as the
required long varying-response witness.
\end{proposition}

\begin{proof}
Relabel the seed site as $0$, put
\begin{equation*}
 t=\frac{1+\alpha}{1-\alpha}>1,
 \qquad \lambda=\frac1{2\rho}>2(t+1),
 \qquad
 y_i=\frac{\vartheta x_{f_i}}{\alpha\rho},
\end{equation*}
and use the normalized coordinates $p_i,s_i,q_i$ from
\cref{eq:double-cycle-feed-normalization}.  Inactive coordinates are zero.
At an active nonseed site the exact equations are
\begin{align}
 t p_i-\frac{p_{i-1}+p_{i+1}+s_i}{4}-\frac{q_i}{2}&=-2,
 &
 t q_i-\frac{p_i}{4}&=-1,
 \notag\\
 t s_i-\frac{s_{i-1}+s_{i+1}+p_i}{4}-\frac{y_i}{2}&=-2,
 &
 t y_i-\frac{s_i}{2}&=-1.
 \label{eq:double-cycle-anchor-active-equations}
\end{align}
The $u_0$ equation has right side $\lambda-2$ instead of $-2$.
The inactive relay, petal, and report gates are respectively
$p_i>4$, $s_i>2$, and $q_i>2$.  On each fixed face all these coordinates are
affine in $\lambda$ with nonnegative slope, by inverse positivity of the
restricted Stieltjes matrix; every coordinate below that is connected to the
seed has strictly positive slope.

The initial face is $\{u_0\}$ and $s_0=(\lambda-2)/t$.  Hence $f_0$ is in
the first batch, while the three other possible violations
$b_0,u_{-1},u_1$ have the same normalized demand $-2+s_0/4$ and are in that
batch exactly when $\lambda>8t+2$.  If only $f_0$ enters, exact settlement on
$\{u_0,f_0\}$ gives
\begin{equation*}
 s_0=\frac{2(2\lambda t-4t-1)}{(2t-1)(2t+1)}.
\end{equation*}
Thus the same triple is the entire next batch exactly when
\begin{equation}
 \lambda>\Lambda_0:=8t+2-\frac{3}{2t};
 \label{eq:double-cycle-anchor-launch}
\end{equation}
at equality the strict gate remains closed, and for
$\lambda\leq\Lambda_0$ the trace stops after $f_0$.  Therefore every trace
that proceeds farther reaches, either in one batch or in two,
\begin{equation*}
 U_{\rm launch}=\{u_0,f_0,b_0,u_{-1},u_1\}.
\end{equation*}
This accounts for the simultaneous initial batches rather than inferring the
chronology from the first petal alone.

Reflection symmetry is preserved by every canonical batch.  At the launch
face it gives $p_0=s_{-1}=s_1=:z$, where
\begin{equation*}
 t s_0-\frac{3z}{4}-\frac{y_0}{2}=\lambda-2,
 \qquad
 t y_0-\frac{s_0}{2}=-1,
 \qquad
 t z-\frac{s_0}{4}=-2.
\end{equation*}
In particular,
\begin{equation*}
 z=\frac{2(2\lambda t-16t^2-4t+3)}{t(16t^2-7)}.
\end{equation*}
The complete set of possible launch-face violations is
\begin{equation}
 \{f_{-1},f_1\}\quad\Longleftrightarrow\quad z>2,
 \qquad
 \{b_{-1},b_1,r_0\}\quad\Longleftrightarrow\quad z>4,
 \qquad
 \{u_{-2},u_2\}\quad\Longleftrightarrow\quad z>8.
 \label{eq:double-cycle-anchor-launch-gates}
\end{equation}
All groups whose displayed strict inequality holds are admitted in the same
all-violations batch.  Later settlement can split the $b_{\pm1}$ and $r_0$
crossings, so we next check their order rather than assuming the launch-face
tie persists.

If $u_2$ becomes positive while inactive, then $f_1$ and $b_1$ are already
positive or co-positive.  The former follows from the thresholds $s_1>8$
and $s_1>2$; for the latter, an inactive $b_1$ has demand
$-2+(p_0+s_1)/4>0$.  The relay $r_0$ is also active or co-positive.  If
$f_1$ is still inactive this is immediate from the launch identity
$p_0=s_1>8$.  Otherwise, before $u_2$ enters, the only two faces with $r_0$
inactive have $b_1$ respectively inactive and active.  At the formal
threshold $s_1=8$, direct solution of
\cref{eq:double-cycle-anchor-active-equations} gives
\begin{equation}
 p_0-4=
 \begin{cases}
 \dfrac{8t^2-3}{2t^2},&b_1\notin U,\\[5pt]
 \dfrac{4(16t^2-5)}{(4t-1)(4t+1)},&b_1\in U,
 \end{cases}
 \label{eq:double-cycle-anchor-precedence}
\end{equation}
which is strictly positive for $t>1$.  Positive face slopes then prove the
claim at every $s_1>8$.

There are no omitted earlier $r_1$ or $b_2$ faces in this comparison.  Once
$f_1,b_1$ are active, the $r_1$ threshold is $p_1=4$.  At that equality,
with $r_0$ respectively inactive and active, the exact values are
\begin{equation*}
 s_1-8=
 \frac{2(32t^3-16t^2-2t+3)}{8t^2-1},
 \qquad
 s_1-8=
 \frac{8(16t^3-8t^2-3t+2)}{16t^2-3},
\end{equation*}
both positive for $t>1$.  Thus $u_2$ is admitted no later than $r_1$;
$b_2$, whose gate before $u_2$ would require $p_1>8$, cannot precede it
either.  The same conclusions hold at the negative indices.

It remains to exhaust the faces on which $f_2$ could first violate.  The
preceding order shows that such a face contains
\begin{equation*}
 \{u_0,f_0,b_0,u_{\pm1},f_{\pm1},b_{\pm1},r_0,u_{\pm2}\}.
\end{equation*}
Before $f_{\pm2}$ or a report enters, its only optional local groups are
$b_{\pm2}$ and $r_{\pm1}$.  Write
$\iota=\mathbf 1_{\{b_{\pm2}\subseteq U\}}$ and
$\delta=\mathbf 1_{\{r_{\pm1}\subseteq U\}}$.  Solve the four corresponding
restricted systems and formally set $s_2=2$.  The resulting report margin
$\Gamma_{\iota\delta}:=(q_0-2)|_{s_2=2}$ is
\begin{equation}
\begin{aligned}
 \Gamma_{00}&=\frac{(2t+1)(8t-5)}{t(16t^2-3)},\\
 \Gamma_{10}&=\frac{(2t+1)(32t^3-16t^2-4t+1)}
 {t(64t^4-8t^2-1)},\\
 \Gamma_{01}&=\frac{64t^4-8t^3-28t^2+t+3}
 {t(2t-1)(2t+1)(4t-1)(4t+1)},\\
 \Gamma_{11}&=\frac{256t^5-128t^3-8t^2+16t+1}
 {64t^4(2t-1)(2t+1)}.
\end{aligned}
 \label{eq:double-cycle-anchor-gamma}
\end{equation}
All four are strictly positive.  For example, after substituting $h=t-1>0$,
the nontrivial numerator polynomials in the last three lines become
\begin{align*}
 &32h^3+80h^2+60h+13,\\
 &64h^4+248h^3+332h^2+177h+32,\\
 &256h^5+1280h^4+2432h^3+2168h^2+896h+137.
\end{align*}
The exact slope ratios $(\partial_\lambda q_0)/(\partial_\lambda s_2)$ on
the same four faces, in the order $00,10,01,11$, are
\begin{equation}
\begin{split}
 &\frac{4(2t-1)(2t+1)}{16t^2-3},\qquad
 \frac{2(2t-1)(2t+1)(8t^2-1)}{64t^4-8t^2-1},\\
 &\frac{128t^4-48t^2+5}
 {2(2t-1)(2t+1)(4t-1)(4t+1)},\qquad
 \frac{512t^6-256t^4+40t^2-1}
 {128t^4(2t-1)(2t+1)}.
\end{split}
 \label{eq:double-cycle-anchor-slope}
\end{equation}
They too are positive for $t>1$.  Hence on every possible face
\begin{equation*}
 s_2>2\quad\Longrightarrow\quad q_0>2.
\end{equation*}
The $f_2$ and $w_0$ violations are therefore admitted in the same canonical
batch, unless $w_0$ entered earlier.

For completeness, no deeper active coordinate can invalidate the four-face
exhaustion.  A new $r_2$ would require $p_2>4$.  On a face with $b_2$ active,
set $p_2=4$ and flag whether $u_2$ and $r_1$ are active.  The four exact
values of $q_0-2$ are
\begin{equation}
\begin{aligned}
 \Delta_{00}&=\frac{(2t+1)(64t^3-16t^2-8t+1)}{t(16t^2-3)},\\
 \Delta_{01}&=\frac{(2t+1)(64t^4-16t^3-16t^2+3t+1)}
 {t^2(16t^2-3)},\\
 \Delta_{10}&=\frac{(2t+1)(64t^3-24t^2-8t+1)}{t(24t^2-5)},\\
 \Delta_{11}&=\frac{(2t-1)(512t^5+320t^4-64t^3-64t^2-8t-1)}
 {t(192t^4-48t^2+1)}.
\end{aligned}
 \label{eq:double-cycle-anchor-deep}
\end{equation}
Every denominator is positive for $t>1$, and substituting $h=t-1$ makes
every displayed numerator polynomial coefficient positive.  Positive slopes
again show that $p_2>4$ forces a report.

The first farther cycle site needs a separate antipodal check when $n=6$.
There $b_3=b_{-3}$ sees both $b_2$ and $b_{-2}$, so its inactive demand is
$-2+p_2/2$ and its strict threshold is $p_2>4$; the
$\Delta_{\iota\delta}$ calculation therefore excludes it before a report.
Likewise $u_3=u_{-3}$ has threshold $s_2>4$, but already
$s_2>2\Longrightarrow q_0>2$ by
\cref{eq:double-cycle-anchor-gamma,eq:double-cycle-anchor-slope}.  For even
$n\geq8$, each first-frontier site $b_{\pm3}$ or $u_{\pm3}$ sees only its one
active $\pm2$ neighbor.  Its respective threshold is $p_2>8$ or $s_2>8$;
the former has already crossed the relay/report threshold, and the latter has
already crossed the petal/report implication.  All still farther vertices
are excluded by the same first-frontier argument.

Thus either the canonical trace stops, a report appears before the next
petal layer, or its first $f_{\pm2}$ batch co-admits $w_0$.  The only petals
in a strictly report-free prefix are consequently
\begin{equation}
 F\cap U\subseteq\{f_0,f_{-1},f_1\},
 \qquad |F\cap U|\leq3.
 \label{eq:double-cycle-anchor-three-petal}
\end{equation}
This is $O(1)$, not $\Omega(n)$.  The required chronology fails before any
response-direction lower bound is needed.
\end{proof}

\begin{proposition}[A relay seed reports in its first propagating batch]
\label{prop:double-cycle-relay-seed-obstruction}
Fix an even $n\geq4$ in the double-cycle graph
\cref{eq:double-cycle-feed-edges}.  For every site $j$, every
$\alpha\in(0,1)$, and every $\rho>0$, run the exact canonical
all-violations gate with relay seed $s=e_{r_j}$.  Put
\begin{equation*}
 t=\frac{1+\alpha}{1-\alpha},\qquad
 \rho_W(\alpha)=\frac{1}{2(2t+1)},\qquad
 \rho_{\mathcal B}(\alpha)=\frac{1}{2(4t+1)}.
\end{equation*}
Then $0<\rho_{\mathcal B}(\alpha)<\rho_W(\alpha)<1/2$, and the
initialized face $U_0$ together with its first exact boundary batch $B_0$
has the exhaustive split below, with $U_0=B_0=\varnothing$ denoting that the
gate does not initialize:
\begin{equation}
 (U_0,B_0)=
 \begin{cases}
  (\varnothing,\varnothing),&\rho\geq1/2,\\
  (\{r_j\},\varnothing),&\rho_W(\alpha)\leq\rho<1/2,\\
  (\{r_j\},\{w_j\}),&\rho_{\mathcal B}(\alpha)\leq\rho<\rho_W(\alpha),\\
  (\{r_j\},\{w_j,b_j\}),&0<\rho<\rho_{\mathcal B}(\alpha).
 \end{cases}
 \label{eq:double-cycle-relay-first-batch}
\end{equation}
At either equality endpoint, the corresponding zero-demand vertex is not
admitted by the strict gate: $w_j\notin B_0$ at
$\rho=\rho_W(\alpha)$, whereas
$b_j\notin B_0$ and $w_j\in B_0$ at
$\rho=\rho_{\mathcal B}(\alpha)$.  Thus every relay-seeded trace that propagates
meets $W$ in its first committed batch, before any petal is admitted.  In
particular this seed orbit supplies no nonempty report-free later-petal
epoch.
\end{proposition}

\begin{proof}
The relay seed has shifted load
\begin{equation*}
 (\ell_\rho)_{r_j}=\frac{\alpha}{\sqrt2}-\alpha\rho\sqrt2.
\end{equation*}
If $\rho\geq1/2$, it and all other shifted loads are nonpositive.  The zero
vector is therefore the unique optimum, including at equality, by strict
convexity and the KKT conditions.  Now suppose $0<\rho<1/2$ and write
$\lambda=1/(2\rho)>1$.  The only positive shifted load is at $r_j$, so
$U_0=\{r_j\}$.  Exact settlement on this singleton gives
\begin{equation}
 x_{r_j}=\frac{\alpha\rho\sqrt2(\lambda-1)}{a},
 \qquad
 q_j:=\frac{\vartheta x_{r_j}}{\alpha\rho\sqrt2}
      =\frac{\lambda-1}{t}.
 \label{eq:double-cycle-relay-seed-coordinate}
\end{equation}
Its only boundary vertices are $w_j$ and $b_j$, whose normalized exact
demands are
\begin{equation}
 \frac{e_{\ell_\rho}(x)_{w_j}}{\alpha\rho\sqrt2}
 =-1+\frac{q_j}{2},
 \qquad
 \frac{e_{\ell_\rho}(x)_{b_j}}{\alpha\rho}
 =-2+\frac{q_j}{2}.
 \label{eq:double-cycle-relay-seed-demands}
\end{equation}
Consequently $w_j$ violates exactly when
$\lambda>2t+1$, equivalently $\rho<\rho_W(\alpha)$, and $b_j$ violates
exactly when $\lambda>4t+1$, equivalently
$\rho<\rho_{\mathcal B}(\alpha)$.  The strict inequalities also give the
two endpoint claims.  Since $\rho_{\mathcal B}<\rho_W$ and there are no
other boundary vertices, these are all possible simultaneous first batches
and prove \cref{eq:double-cycle-relay-first-batch}.  The calculation is
local to site $j$ and does not use the matching neighbor of $w_j$; hence it
covers every rotation and both parity classes of the fixed matching on $W$.
\end{proof}

\begin{remark}[Seed and batching scope, and the full double-cycle ledger]
\label{rem:double-cycle-feed-trace-ledger}
The backbone-seed restriction is substantive.  If the seed is instead
$u_j$ and
\begin{equation}
 0<\rho<\frac{1-\alpha}{8},
 \label{eq:double-cycle-anchor-seed-counterrange}
\end{equation}
then the initial face is $\{u_j\}$, with
$x_{u_j}=\alpha(1/2-2\rho)/a$, and
\begin{equation*}
 e_{\ell_\rho}(x)_{f_j}
 =-\alpha\rho+\frac{\vartheta}{2}x_{u_j}>0,
 \qquad
 e_{\ell_\rho}(x)_{w_i}=-\alpha\rho\sqrt2<0
 \quad(1\leq i\leq n).
\end{equation*}
Thus an anchor seed gives an exact report-free first-petal batch; no
all-seed-orbit version of \cref{prop:double-cycle-feed-report-obstruction}
follows.  The complete canonical continuation is nevertheless stopped by
\cref{prop:double-cycle-anchor-seed-obstruction}: at most the three petals at
sites $j,j-1,j+1$ precede the first report, and the first $j\pm2$ petal batch
co-admits $w_j$.  A report seed violates
report exclusion initially, a petal seed is an initial attachment rather than
a later admission.  The remaining relay-seed orbit is exhausted by
\cref{prop:double-cycle-relay-seed-obstruction}: it either has no initialized
face, terminates on its singleton face, or admits $w_j$ in its first batch,
possibly together with $b_j$, and admits no petal first.  Thus the backbone,
anchor, and relay candidate orbits all fail the prescribed long report-free
later-petal trace.  This retires this double-cycle graph family for that
witness, not all possible cyclic witnesses.  The proofs use canonical all-violations
batching.  A positive-subset policy could defer a simultaneously positive
report and is not covered.

These are exact KKT zero-sign statements at regularization level $\rho$, not
finite-band, word/bit, finite-precision, or stability claims.  Although
$Q_{\mathcal B A}=-(\vartheta/4)I$,
$Q_{\mathcal B R}=-(\vartheta/\sqrt8)I$,
$Q_{AF}=-(\vartheta/2)I$, and
$Q_{RW}=-(\vartheta/2)I$ have full rank, the failed canonical chronology
shows only that these cuts do not realize the prescribed report-free
petal/attachment trace with changing response directions.  It proves nothing
about response directions after $W$ enters and yields neither a reporter,
response-rank, nor work lower bound.

Use the shared order
\begin{equation*}
 \mathbf L=(C_{\rm adj},R_{\rm adj},R_{\rm int},C_{\rm pre},C_{\rm ctl},
 C_{\rm rec},C_{\rm resp},M_{\rm pers},M_{\rm tmp},C_{\rm mat},C_{\rm emit}).
\end{equation*}
For any tested canonical prefix, let $\mathbf P_k$ be its complete paid
eleven-vector: it includes graph exposure and repeated row reads, adjacency
and external rounds, preprocessing, all exact-face recurrence or response
work, gate/control work, retained and scratch state, recovery/materialization,
and every task-required emission through face $k$.  The graph has
$\operatorname{cvol}(V)=18n$.  If its bounded-degree rows have not been
retained, one nonadaptive streamed validation of all $n$ implications above
has incremental vector
\begin{equation}
 \Delta\mathbf L_{\rm fresh}
 =\bigl(\Theta(n),1,0,0,\Theta(n),0,0,0,O(1),0,0\bigr).
 \label{eq:double-cycle-feed-fresh-audit-vector}
\end{equation}
If those rows were already retained and charged in
$M_{\rm pers}(\mathbf P_k)$, validation has
\begin{equation}
\Delta\mathbf L_{\rm stored}
 =\bigl(0,0,0,0,\Theta(n),0,0,0,O(1),0,0\bigr).
 \label{eq:double-cycle-feed-stored-audit-vector}
\end{equation}
The total ledger is $\mathbf P_k\oplus\Delta\mathbf L$, where additive
coordinates are summed and the two memory coordinates take their larger
peak.  Constructing a cache during validation is the fresh case with its
writes in $C_{\rm ctl}$ and retained cells in $M_{\rm pers}$, not free reuse.
If batch labels are required, their writes belong to $C_{\rm emit}$; the
mathematical proposition itself requires no intermediate output.  All
numerical face solves, old-face reads, response queries, state writes, and
materialization remain in $\mathbf P_k$ and are not hidden by either
incremental validation vector.

For the fixed anchor seed in
\cref{prop:double-cycle-anchor-seed-obstruction}, let
$\mathbf P_{\rm firstW}$ be the complete paid prefix through the first report
batch, or through termination if there is no report.  The local exact
chronology and its constant-radius case tables have incremental vectors
\begin{equation}
 \Delta\mathbf L^{\rm anchor}_{\rm fresh}
 =\bigl(O(1),1,0,0,O(1),0,0,0,O(1),0,0\bigr),
 \qquad
 \Delta\mathbf L^{\rm anchor}_{\rm stored}
 =\bigl(0,0,0,0,O(1),0,0,0,O(1),0,0\bigr).
 \label{eq:double-cycle-anchor-audit-vectors}
\end{equation}
Auditing every rotated anchor site instead restores the $\Theta(n)$ fresh
scan/control or stored control entries above; a streamed scratch peak remains
$O(1)$.  In either form the total is the fully charged prefix
$\mathbf P_{\rm firstW}\oplus\Delta\mathbf L^{\rm anchor}$; no face solve,
batch emission, or response work is erased by the incremental proof audit.

For the fixed relay seed in
\cref{prop:double-cycle-relay-seed-obstruction}, let
$\mathbf P^{\rm relay}_{\rm firstW}$ be the complete paid prefix through the
first report batch, or through termination or the absent initialization in
the other two cases of \cref{eq:double-cycle-relay-first-batch}.  Only the
constant-radius proof audit is incremental:
\begin{equation}
 \Delta\mathbf L^{\rm relay}_{\rm fresh}
 =\bigl(O(1),1,0,0,O(1),0,0,0,O(1),0,0\bigr),
 \qquad
 \Delta\mathbf L^{\rm relay}_{\rm stored}
 =\bigl(0,0,0,0,O(1),0,0,0,O(1),0,0\bigr).
 \label{eq:double-cycle-relay-audit-vectors}
\end{equation}
The total is
$\mathbf P^{\rm relay}_{\rm firstW}\oplus
\Delta\mathbf L^{\rm relay}$, so every exact-face solve, row read, gate
decision, state write, and task-required batch emission remains charged in
the prefix; there is no reporter or response work hidden in the local vector.
Auditing every rotated relay site restores the $\Theta(n)$ fresh scan/control
or stored control entries of the global audit, while scratch remains $O(1)$.
All of these relay statements are exact-real canonical claims only, with no
arbitrary positive-subset, post-report, reporter/work lower-bound,
finite-precision, or stability consequence.
\end{remark}

\begin{problem}[Finite-resolution cyclic response]
\label{prob:finite-resolution-response}
For $\eta=\alpha\tau$, maintain the certificates
\cref{eq:finite-response-certificate}, or the strictly weaker interval
certificates \cref{eq:finite-response-intervals}, and report the
finite-resolution gate
with
\begin{equation}
 \mathcal C_{\rm resp}^{(\eta)}
 =\widetilde O\!\left(
  \frac{1+\operatorname{cvol}(S^\star(\rho))}{\sqrt\alpha}
 \right).
 \label{eq:finite-resolution-response-target}
\end{equation}
No exact positive-versus-nonpositive distinction is required inside the
normalized KKT band $[0,\eta]$.  By
\cref{lem:finite-response-energy-packing}, the degree volume whose cumulative
normalized response moves through a prescribed interval margin is charged to
objective shock.  The remaining operation is therefore an output-sensitive
heavy-change report, not an eager refresh of every boundary demand.  In the
certified response-tree model, the exact theorem supplies this oracle only
conditionally on bounding visited nodes, finite-accuracy factor maintenance,
leaf work, and presentation updates at the displayed scale.
\end{problem}

\Cref{prop:notched-sun-rppr-trace-obstruction} rules out the existing
notched-double-sun structural epoch as a witness for this problem: its source
and report leaves, apart from the at-most-one pair containing the seed, cross
the exact gate together in matched pairs.  The problem remains open on legal
cyclic traces whose source and report rows are genuinely asymmetric.
\Cref{prop:matched-report-sun-gate-obstruction} tests a minimal edge-only
degree lift of those rows and rules it out too.  The matching makes the raw
report cut full rank, but report nonactivation imposes an anchor-coordinate
upper bound incompatible with a nonseed petal violation.  A next candidate
must therefore attenuate the report through an additional active stem or a
different bounded-degree threshold gadget; merely raising the degree of a
direct report neighbor cannot supply the desired epoch.
\Cref{prop:active-relay-matched-sun-preload-obstruction} performs that first
stem audit and stops it one step earlier: the internal degree-two relay cannot
become active before its degree-one petal.  The canonical all-violations gate preserves
this precedence in canonical all-violations batches of arbitrary size and
under every seed placement, with
only the one relay-seed site excepted.  Thus the next graph must reverse the
petal/relay thresholds or feed the relay from an independent active backbone;
there is still no legal cyclic trace on which reporter algebra may begin.
\Cref{prop:petal-cycle-relay-first-obstruction} tests the literal threshold
reversal and rules out its intended chronology as well.  Although the
degree-three petal has the isolated threshold $6\alpha\rho/\vartheta$, a
report-quiet all-relay face already contains all but at most the petal paired
with an anchor or relay seed.  The obstruction is the coupled anchor balance,
not a failure of the local threshold window or a response-rank theorem.  An
independently active relay backbone, or a different gadget whose global
settlement avoids this cap contradiction, is therefore the next legality
test.
\Cref{prop:double-cycle-feed-report-obstruction} performs the literal
independent-backbone test and stops it more strongly for every backbone seed:
as long as the next all-violations batch remains report-free, every active
anchor stays at or below its petal threshold and every newly admitted anchor
finds its relay already active or co-admits it in the same batch.  Thus no
petal enters before a report.  The full-rank raw cuts do not realize the
prescribed report-free petal/attachment sequence with changing report
responses; nothing is claimed after $W$ enters.  The proof is not
seed-uniform---\cref{eq:double-cycle-anchor-seed-counterrange} gives an exact
anchor-seed first-petal counterexecution---but
\cref{prop:double-cycle-anchor-seed-obstruction} follows that counterexecution
to the first report and permits only three report-free petals.
\Cref{prop:double-cycle-relay-seed-obstruction} exhausts the remaining core
orbit: its exact singleton demands have thresholds
$\rho_W=1/(2(2t+1))$ and
$\rho_{\mathcal B}=1/(2(4t+1))$, so the relay face either terminates or its
first propagating batch contains $w_j$, possibly together with $b_j$.
Accordingly the present double-cycle graph family is retired for the
prescribed long report-free later-petal witness.  These results do not cover
positive-subset batching, another cyclic family, or any post-report trace.
A next candidate must therefore decouple the report cap from the anchor
maximum principle, rather than merely duplicating the anchor cycle as a relay
feed.

This problem is strictly weaker than exact light-sign maintenance.  If it is
proved, then for $0<\varepsilon\leq1$, choosing
$\rho=\tau_{\rm gate}=\tau_{\rm sol}=\varepsilon/3$ in
\cref{thm:finite-resolution-continuation} gives
\begin{equation}
 \|D^{-1/2}(x^0-z)\|_\infty\leq\varepsilon,
 \qquad
 \Work=\widetilde O\!\left(\frac1{\varepsilon\sqrt\alpha}\right).
 \label{eq:finite-resolution-product}
\end{equation}
This is a note-scoped degree-normalized PageRank guarantee; identifying it
with the repository's final residual convention still requires the explicit
accuracy mapping excluded in \cref{sec:scope}.

A stronger local sufficient primitive queues a consecutive light group at
anchor face $U$ using one old-face response scan and work proportional to the new charged
volume.  The lazy fan and kinetic spider constructions already have this
property.  Proving it for arbitrary nonequitable thick cores is the remaining
graph-uniform question.  Together with \cref{thm:aggregate-repair-oracle}, it
gives
\begin{equation}
 \widetilde O(\operatorname{cvol}(S^\star)/\sqrt\alpha)
 =\widetilde O(1/(\rho\sqrt\alpha))
 \label{eq:aggregate-product-consequence}
\end{equation}
total work.
